\abstract{%
Quantitative research is an iterative loop: a researcher proposes a market hypothesis, implements it as an executable signal, tests it on historical data, and decides what to keep. Existing benchmarks score formulas, code, or individual trading decisions, leaving open whether large language models (LLMs) can complete this loop and produce strategies that work beyond the research period. We introduce \textbf{\BenchmarkName}, a benchmark and public leaderboard for end-to-end quantitative research on real A-share market data. A model constructs features and stock-scoring code, receives feedback only on 2022--2023, and freezes one strategy for assessment on the sealed year 2024. A cash-account evaluator applies trading constraints and costs, including market impact estimated from a simulated order book. Scores combine signal quality, excess return, token efficiency, and the diversity of features used by the strategy. We evaluate six frontier models in one-shot and iterative settings. Iteration raises 2024 returns for five of the six models. No one-shot strategy beats the equal-weight market, whereas the best iterative strategy (\ModelClaudeOpus) earns 13.3 percentage points of excess return. In a two-model feedback ablation, withholding performance metrics while retaining execution errors lowers mean 2024 returns for both models. Higher predictive correlation can still accompany lower portfolio returns, and all evaluated models trail Kepler-1, our closed-source professional strategy. \BenchmarkName\ is open to submissions from any model or agent harness.
}
